\chapter{Методические указания}

В разделе <<Методические указания>> приведена программа курса Статика корабля, соответствующие ей вопросы для самопроверки, а также список литературы.

\section{Программа курса Статика корабля}

\textit{Тема 1. Основные положения}

Статика корабля как один из разделов теории корабля.
Краткая история развития учения о плавучести и остойчивости.
Теоретический чертеж судна.
Главные размерения, их соотношения, коэффициенты полноты, системы координат.

\vspace{0.5cm}

Вопросы и задания для самопроверки:

1.	Что представляет собой теоретический чертеж судна

2.	Нарисуйте систему координат, принятую в статике корабля.

\vspace{0.5cm}

\textit{Тема 2. Плавучесть}

Посадка судна и ее параметры.
Силы, действующие на плавающее судно.
Условия равновесия судна.
Объемное водоизмещение и координаты центра величины при посадке судна прямо и на ровный киль.
Строевые по ватерлиниям и шпангоутам.
Кривая водоизмещения и грузовой размер.
Элементы площади ватерлиний.
Кривые абсцисс центра величины и центра тяжести площадей ватерлиний.
Кривая аппликат центра величины.
Площадь шпангоута. Интегральные кривые В.Г.~Власова.
Масштаб Бонжана.
Определение водоизмещения и координат центра величины при посадке судна прямо и с дифферентом.
Диаграмма Г.А.~Фирсова.
Кривые элементов теоретического чертежа.
Запас плавучести.
Грузовая марка.
Условия приема и расходования малого и большого груза без крена и дифферента.
Влияние изменения плотности воды на посадку судна.

\vspace{0.5cm}

Вопросы и задания для самопроверки:

1.	Какими параметрами характеризуется посадка судна? Какими бывают частные случаи посадки судна?

2.	Дайте определение центра тяжести и центра величины судна. Какие силы действуют на плавающее судно? 
Условия равновесия судна.

3.	Напишите формулы для объемного водоизмещения и координат центра величины при посадке прямо и на ровный киль.

4.	Укажите свойства строевых по ватерлиниям и шпангоутам.

5.	Будут ли равны площади, ограниченные строевой по шпангоутам и строевой по ватерлиниям и соответствующими осями координат?

6.	Какая связь существует между строевой по ватерлиниям и грузовым размером?

7.	Напишите формулы для площадей ватерлинии, абсциссы ее центра тяжести и моментов инерции относительно продольной и поперечной осей, проходящих через ее центр тяжести.

8.	Какая зависимость существует между кривой абсцисс центра величины и кривой центра тяжести площади ватерлиний?

9.	Какие задачи решаются с помощью интегральных кривых площадей шпангоутов В.Г. Власова?

10.	Объясните принцип построения масштаба Бонжана и укажите, какие задачи решаются с его помощью.

11.	Объясните методику построения диаграммы Г.Н.~Фирсова.

12.	Какие зависимости включают в себя кривые элементов теоретического чертежа?

13.	Что такое запас плавучести и какими мерами он обеспечивается?

14.	Объясните назначение грузовой марки.

15.	Сформулируйте условия приема на судно малого и большого грузов без крена и дифферента.

16.	Как влияет изменение плотности воды на посадку судна?

\vspace{0.5cm}

\textit{Тема 3. Остойчивость}

Общее понятие об остойчивости.
Теорема Эйлера для равнообъемных ватерлиний.
Перемещение ЦВ при малых равнообъемных наклонениях.
Метацентры и метацентрические радиусы.
Восстанавливающий момент.
Плечо статической остойчивости.
Начальная остойчивость.
Понятие о метацентрической высоте.
Расчет плеч статической остойчивости при равнообъемных наклонениях.
Интерполяционные кривые плеч остойчивости формы.
Диаграмма статической остойчивости.
Работа восстанавливающего момента.
Диаграммы динамической остойчивости.
Решение задач с помощью диаграмм статической и динамической остойчивости.

\vspace{0.5cm}

Вопросы и задания для самопроверки:

1.	Дайте определение остойчивости корабля.

2.	Сформулируйте теорему Эйлера для равнообъемных наклонений.

3.	Покажите, как перемещается центр величины при равнообъемных наклонениях.

4.	Что такое метацентры и метацентрические радиусы, как вычисляются метацентрические радиусы?

5.	Разъясните природу возникновения восстанавливающего момента при накренении судна.
Что такое плечо статической остойчивости?

6.	Что является мерой начальной остойчивости?
Напишите формулы для начальной метацентрической высоты.

7.	Как определяются плечи остойчивости формы и остойчивости веса?

8.	Объясните методику построения и назначение интерполяционных кривых плеч остойчивости формы.

9.	Изобразите диаграмму статической остойчивости с положительной и отрицательной начальными метацентрическими высотами.

10.	Что такое динамическая остойчивости судна и что представляет собой диаграмма динамической остойчивости?

11.	Укажите связь между плечами динамической и статической остойчивости и обобщенной метацентрической высотой.

12.	Какая связь существует между диаграммами статической и динамической остойчивости?

13.	Укажите все характерные точки на диаграмме статической остойчивости, определите обобщенную метацентрическую высоту при произвольном угле крена.

14.	Как определить угол крена судна при воздействии на него произвольно заданного кренящего момента, приложенного статически?

15.	Как определить предельный угол крена судна и предельный момент, приложенный статически?

16.	Как по диаграмме статической остойчивости определяется угол крена от динамически приложенного кренящего момента?

17.	Как определить по диаграмме статической и динамической остойчивости предельный угол крена и предельный кренящий динамически приложенный момент?

18.	Объясните методику и последовательность расчета плеч остойчивости по второму способу Крылова–Дарньи.

\vspace{0.5cm}

\textit{Тема 4. Практические приложения теории плавучести и остойчивости}

Изменение посадки и остойчивости судна при переносе, приеме и расходовании груза.
Понятие о нейтральной плоскости.
Влияние на остойчивость подвешенного, перекатывающегося, жидкого и сыпучего груза.
Изменение остойчивости судна при движении на волнении.
Нормирование остойчивости.
Основные требования Правил Российского морского регистра судоходства.
Определение положения ЦТ судна опытным путем.
Методика проведения опыта кренования и обработка его результатов.

\vspace{0.5cm}

Вопросы и задания для самопроверки:

1.	Как меняется посадка и остойчивость судна при вертикальном и горизонтальном переносах груза?

2.	Поясните методику определения изменения остойчивости судна при приеме и расходовании малого груза.

3.	Дайте определение нейтральной плоскости, напишите ее уравнение.

4.	В какой последовательности решается задача об определении изменения остойчивости и посадки судна при приеме и расходовании большого груза?

5.	как оценить влияние подвешенного груза на остойчивость судна?

6.	В каких случаях, почему и как жидкий груз может оказать влияние на остойчивость судна?

7.	Как уменьшить вредное влияние жидкого груза на остойчивость судна?

8.	Разъясните физическую природу процесса влияния сыпучего груза на остойчивость судна.

9.	Какое положение судна относительно набегающих волн наиболее опасно с точки зрения потери остойчивости и почему?

10.	Что такое критерий погоды и как он определяется?

11.	Какие требования предъявляют Правила Российского морского регистра судоходства к диаграммам статической остойчивости и начальной метацентрической высоте судов?

12.	Опишите методику выполнения опыта кренования и объясните цель его проведения.

\vspace{0.5cm}

\textit{Тема 5. Непотопляемость}

Основные определения и требования.
Типы затопленных отсеков.
Два способа расчета непотопляемости.
Учет прямостенности борта в расчетах непотопляемости судна.
Расчет непотопляемости при затоплении группы отсеков.
Диаграмма непотопляемости С.Н. Благовещенского.
Кривые предельных длин отсеков.
Практические способы спрямления аварийного судна.
Требования Российского морского регистра судостроения к элементам непотопляемости транспортных судов.

\vspace{0.5cm}

Вопросы и задания для самопроверки:

1.	Дайте определение непотопляемости.

2.	Разъясните, какой тип затапливаемых отсеков наиболее опасен и почему.

3.	Какие существуют способы расчета непотопляемости и для каких типов затапливаемых отсеков они используются?

4.	Объясните подробно суть способа постоянного водоизмещения.

5.	Как соотносятся между собой значения метацентрических высот аварийного судна, определенные способами приема груза и постоянного водоизмещения?

6.	Опишите схему определения осадки аварийного судна с учетом прямостенности борта способами приема груза и постоянного водоизмещения.

7.	Что понимается под определением «эквивалентных отсек» при затоплении группы отсеков?

8.	Нарисуйте диаграмму С.Н. Благовещенского и разъясните методику ее использования.

9.	Дайте определение предельной линии погружения, предельной аварийной ватерлинии и предельной длины отсека.

10.	Расскажите последовательность построения кривой предельных длин отсеков.

11.	Опишите пять характерных случаев состояния аварийного судна.
Какие методы спрямления используются на практике?

12.	Что представляют собой таблицы непотопляемости А.Н. Крылова?

13.	Какие требования предъявляет Российский морской регистр судоходства к элементам остойчивости поврежденного судна?

\vspace{0.5cm}

\textit{Тема 6. Спуск судов на воду}

Основные понятия и определения.
Способы спуска судов. Основные элементы и характеристики спускового устройства при продольном спуске.
Силы, действующие на судно при спуске.
Деление продольного спуска на периоды.
Первый период спуска.
Второй период спуска.
Опрокидывание и всплытие судна при спуске.
Третий и четвертый периоды спуска.
Диаграмма продольного спуска.
Боковой спуск.

\vspace{0.5cm}

Вопросы и задания для самопроверки:

1.	Какие виды спуска судна существуют?

2.	Опишите основные элементы спускового устройства.

3.	Какие силы действуют на судно при спуске до входа в воду и после его входа в воду?

4.	Что такое потерянная сила плавучести?

5.	Кратко охарактеризуйте периоды продольного спуска.

6.	Может ли произойти опрокидывание судна при продольном спуске, если его центр тяжести не перешел порог?

7.	Какие меры принимаются для предотвращения опрокидывания?

8.	Как определяется момент начала всплытия при спуске?

9.	Что такое баксово давление?

10.	Разъясните явление соскакивания при спуске, вероятность его появления, возможные последствия.

11.	Расскажите, как определяются возможности опрокидывания, момент начала всплытия и величина баксова давления по английской и французской диаграммам спуска.

12.	В чем состоит принципиальное отличие английской и французской диаграмм спуска?

13.	Объясните различие между явлениями опрокидывания при продольном и боковом спусках?


\section{Литература}

\begin{enumerate}
    \item Борисов~Р.В. Статика корабля: учебник / Р.В.~Борисов. --- Москва ; Вологда : Инфра-Инженерия, 2023. --- 216 с. : ил., табл. --- ISBN 978-5-9729-1126-4
    \item Борисов~Р.В. Статика корабля: учебное пособие для вузов / Р.В.~Борисов, В.В.~Луговской, В.В~Рождественский, Б.В.~Мирохин~Б.В. --- 2-е изд., перераб. и доп. --- СПб.: Судостроение, 2005. --- 256 с. --- ISBN 5-7355-0634-X
    \item Благовещенский~С.Н.Справочник по статике и динамике корабля: в 2 томах / С.Н.~Благовещенский, А.Н.~Холодилин. \\
    Том 1. Статика корабля. --- Л.: Судостроение, 1975. --- 336 с.
    \item Мирохин~Б.В.Теория корабля / Б.В.~Мирохин, Б.В.~Жинкин, Г.И.~Зильман --- Л.: Судостроение, 1989. --- 352 с.
    \item Справочник по теории корабля: В 3 томах / Под ред. Я.И.~Войткунского \\
    Том 2. Статика судов. Качка судов / Под ред. Я.И.~Войткунского --- Л.: Судостроение, 1985. --- 440 с.
\end{enumerate}